% Generated draft: requires mathematical and source review.
\documentclass[11pt]{article}
\usepackage[margin=25mm]{geometry}
\usepackage{fontspec}
\setmainfont{Latin Modern Roman}
\usepackage{amsmath,amssymb,mathtools,unicode-math}
\setmathfont{Latin Modern Math}
\usepackage{enumitem,needspace,xurl,hyperref}
\setlength{\parindent}{0pt}
\setlength{\parskip}{5pt}
\setlength{\emergencystretch}{3em}
\newcommand{\N}{\mathbb{N}}
\newcommand{\Z}{\mathbb{Z}}
\newcommand{\Q}{\mathbb{Q}}
\newcommand{\R}{\mathbb{R}}
\newcommand{\C}{\mathbb{C}}
\newcommand{\F}{\mathbb{F}}
\newcommand{\A}{\mathbb{A}}
\newcommand{\PP}{\mathbb{P}}
\newcommand{\E}{\mathbb{E}}
\newcommand{\eps}{\varepsilon}
\newcommand{\dd}{\,\mathrm{d}}
\begin{document}
\section*{UnsolvedMath ID: 10400178}
\textbf{Source problem code:} catalogue entry
\subsection{Definitions and mathematical statement}
Let \(\lambda_n\) denote the Cappell--Lee--Miller Casson-type \(\mathrm{SU}(n)\) invariant meant by the source. Its definition, manifold domain, admissible ranks and normalization are not given in this extract. C. Lescop's Question 10.3 asks whether these invariants are of finite type and, if so, what their degrees and weight systems are.

To make the finite-type condition precise one must choose the source's surgery or local-move filtration. In general, let \(\mathcal H\) be the intended manifold class, \(R\) the coefficient ring, and
\[
 R\mathcal H=F_0\supset F_1\supset F_2\supset\cdots
\]
the chosen filtration on the free module generated by its equivalence classes. After linear extension, \(\lambda_n\) is of type at most \(d\) if
\[
 \lambda_n(F_{d+1})=0.
\]
Its exact degree requires the appropriate minimal such bound and conventions for constant or zero invariants. At degree \(d\), the induced functional on \(F_d/F_{d+1}\) is its leading symbol; a weight-system description additionally requires the source's identification of this graded quotient with a diagram space and its relations.

The problem asks for these determinations for the given family, with degree and weight-system requests conditional on finite typeness. No particular surgery filtration, universal degree formula, diagram space, or rank restriction is supplied, so none is inserted here. The supplied partial-progress assessment says the answer is known in some cases but does not identify those cases or cite proofs.
\subsection{Short English statement}
Determine whether the specified Casson-type \(\mathrm{SU}(n)\) invariants are finite type. For those that are, find their exact degrees and leading weight systems. The invariant definitions, rank range and finite-type filtration are absent from the source extract.
\subsection{Sources}
\begin{itemize}
\item Local UnsolvedMath registry, permanent ID \(10400178\), exact source code \texttt{catalogue entry}, Question 10.3, attributed to C. Lescop. The title is truncated after ``(C.''
\item Ohtsuki, \emph{Problems on invariants of knots and 3-manifolds} (2002), Question 10.3, PDF page 143: \url{https://msp.org/gtm/2002/04/gtm-2002-04-024s.pdf}. Bold headings distinguish 231 numbered items from references in surrounding exposition.
\item Extraction pdf-font-structure, accessed 2026-07-29. Source evidence NEEDS\_REVIEW presents an open problem and calls for a current primary-literature audit. Registry status partially\_solved and classification PARTIAL-PROGRESS report known results in some cases, with no complete citations or stated cases. No separate reference defining the Cappell--Lee--Miller invariants is supplied.
\item Dataset attribution: Ulam AI / UnsolvedMath Contributors, CC BY 4.0. Source rights NEEDS\_REVIEW: publicly accessible author/publisher copy, redistribution terms requiring review.
\end{itemize}
\end{document}
